\documentclass[10pt,a4paper]{amsart}
\usepackage[margin=19mm]{geometry}
\usepackage{amsmath,amssymb,mathtools}
\usepackage[scr=rsfs,cal=euler]{mathalfa}
\usepackage{xfrac}
\usepackage[unicode,hidelinks,bookmarks=false]{hyperref}
\newcommand{\ie}{i.e.\ }
\newcommand{\cf}{cf.\ }
\newcommand{\resp}{respectively\ }
\newcommand{\Subsection}{\subsection}
\providecommand{\nobreakdash}{\nobreak\hbox{-}}
\setlength{\emergencystretch}{2em}
\multlinegap0pt
\usepackage{dsfont}
\usepackage{amssymb}
\usepackage{xcolor}
\usepackage{breqn}
\usepackage[shortlabels]{enumitem}
\newtheorem{proposition}{Proposition}
\newcommand{\Half}{\R \times \R_+}
\newcommand{\Leb}{\mathrm{Leb}}
\newcommand{\R}{\mathbb{R}}
\newcommand{\e}{\mathbb{E}}
\newcommand{\landsc}{\mathrm{L}}
\newcommand{\zone}{\mathrm{Z}}
\title{The Relay Random Tree:\\A stochastic geometry approach of multihop relay in an urban visibility setting.}
\author{Paul Rax, Fran\c{c}ois Baccelli}
\date{Source: arXiv:2605.03851v1 (CC BY 4.0)}
\begin{document}
\maketitle

\noindent\emph{Source and licence.} The Relay Random Tree:\\A stochastic geometry approach of multihop relay in an urban visibility setting.. Version arXiv:2605.03851v1 (CC BY 4.0), distributed by arXiv under the Creative Commons Attribution 4.0 International licence, http://creativecommons.org/licenses/by/4.0/, as recorded in the arXiv licence metadata for this exact version and shown on the abstract page https://arxiv.org/abs/2605.03851v1. Full licence notice: \texttt{licenses/arxiv-2605.03851-CC-BY-4.0.txt}.\par\smallskip

\noindent\textit{Editorial scope.} Counted unit: the article's proposition mass\_transport\_save (source0.tex lines 925--932). The article's complete original proof of it is delivered with it (source0.tex lines 934--980, one \texttt{proof} environment of 47 lines). No mathematics is written, repaired or re-derived: the statement and the proof are the article's own lines, in the article's own notation, and no retained line is substituted or re-flowed.  No further source-local context is needed: every cross-reference of every retained line resolves inside the two delivered spans.  Nothing was omitted from any retained span, so the package carries no omission record.  No retained line contains a citation, so the wrapper carries no bibliography and no bibliography entry.  No retained line contains a figure environment, an image-inclusion command or drawing code, so no image is delivered and the package declares no asset dependency.  Editorial formatting only, in addition: an AMS \texttt{amsart} preamble that restates the article's own declarations and macros used by the retained lines, namely the environments \texttt{proposition} on separate counters, together with the standard packages those lines need.  The retained environments print in this document as Proposition 1 in order of appearance, whereas the article prints the same items with the numbering of its own full document.  The complete native source of the article as distributed in the arXiv e-print for this exact version is bundled unchanged as \texttt{sources/199822/source0.tex}.
\begin{proposition}\label{mass_transport_save}\textbf{Length of the typical relaying zone.}

Let us consider $H$ a random variable distributed following $\mathcal{L}$. Let us further consider the landscape $\tilde{\landsc} := \landsc\cup \{(0, H)\}$. Then expectation of the length of the zone of the typical eventual relaying $\zone (0, H)$ is 
\[\e[\ell(\zone(0, H))] = \frac{1}{\lambda}\e[T],\]
where $T$ is the hitting time of the cemetery state $(-\infty, -\infty, -\infty)$ for the Markov chain $(\tilde{x}_n, \mathrm{t}_n, H_n)$ started at 
$(0, 0, +\infty)$:
\[\e[T] = \sum_{n\geq 1} \displaystyle\int_{\Half^{n+1}} g_n \delta_{(0, 0)}\displaystyle\bigotimes_{i=1}^{n} (\Leb \otimes \mathcal{L}).\]
\end{proposition}

\begin{proof}
This comes from a simple mass transport principle argument: Let us consider $\Phi = \{x_i\}$ a Poisson point process of intensity $\lambda'>0$ on $\R$ independent from $\tilde{\landsc}$. Then, by definition,
\begin{align*}
\e[\ell(\zone(0, H))] &= \frac{1}{\lambda'}\e[\Phi\cap \zone(0, H)]\\
&= \frac{1}{\lambda'}\e[\# \{x\in \Phi: (x, 0)\text{ is eventually relayed by } (0, H)\}].
\end{align*}
However, the superposition of PPPs allows us to evaluate this quantity.

Consider the landscapes $\landsc_1$, $\tilde{\landsc}_1$ and $\tilde{\landsc}_2$ defined as follows:
\begin{align*}
    &\tilde{\landsc}_1 := \tilde{\landsc}\cup \{(x, 0): x\in \Phi\}\\
    &\tilde{\landsc}_2 := \landsc\cup \{(x, 0): x\in \Phi\}\cup \{(0, 0)\}.
\end{align*}
Clearly, $\landsc_1$ is a Poisson building process of intensity $\lambda +\lambda'$ and with height distribution
\[\frac{\lambda}{\lambda+\lambda'}\mathcal{L}+\frac{\lambda'}{\lambda + \lambda'} \delta_0.\]
Where by abuse of notation, we consider buildings with height $0$. By convention, we consider no relaying by such buildings: we consider the extension of the scheme to be
\[\tau (x, y, \landsc) = (x_b, y_b) := \underset{\substack{(a, b) \in \landsc \\a\in ]x, x+R]\\b>0}}{\text{argmax}} \frac{b-y}{a-x}.\]
As such, if $\omega$ denotes a Bernoulli random variable of parameter $\frac{\lambda}{\lambda+\lambda'}$ independent from everything we defined so far, the landscape
$$\tilde{\landsc'}:=\left\{
    \begin{array}{ll}
        \tilde{\landsc}_1 & \mbox{if } \omega = 1 \\
        \tilde{\landsc}_2 & \mbox{if }\omega=0
    \end{array}
\right.$$
follows the Palm measure of the Poisson Building process $\landsc'$. As such, the relay forest is unimodular when rooted in $(0, \omega H)$. Moreover, for any real function $f$ covariant with rooted graph isomorphisms
\begin{align*}
    \e[f(\mathcal{G}( \tilde{\landsc}'))]  &= \frac{\lambda}{\lambda+\lambda'}\e[f(\mathcal{G}( \tilde{\landsc}_1))] + \frac{\lambda'}{\lambda+\lambda'}\e[f(\mathcal{G}( \tilde{\landsc}_2))].
\end{align*}
Using the mass transport principle, we get on the one hand:
\begin{align*}
    \e[\ell(\zone(0,\omega H, \tilde{\landsc}'))] &= \e\left[\sum_{(x, h)\in \tilde{\landsc}'} \mathds{1}_{h = 0, (x, h) \text{ is eventually relayed by } (0, \omega H)}\right]\\
    &= \frac{\lambda}{\lambda+\lambda'}\e\left[\sum_{(x, h)\in \tilde{\landsc}_1} \mathds{1}_{h = 0, (x, h) \text{ is eventually relayed by } (0,  H)}\right]+ 0
\end{align*}
and on the other hand:
\begin{align*}
    \e\left[\sum_{(x, h)\in \tilde{\landsc}'} \mathds{1}_{h = 0, (x, h) \text{ is eventually relayed by } (0, \omega H)}\right]&= \e\left[\sum_{(x, h)\in \tilde{\landsc}'} \mathds{1}_{\omega H= 0, (0, 0) \text{ is eventually relayed by } (x, h)}\right]\\
    &= 0+ \frac{\lambda'}{\lambda+\lambda'}\e\left[\sum_{(x, h)\in \tilde{\landsc}_2} \mathds{1}_{(0, 0) \text{ is eventually relayed by } (x,  h)}\right].
\end{align*}
However, under our scheme, we can rewrite:
\[\e\left[\sum_{(x, h)\in \tilde{\landsc}_2} \mathds{1}_{(0, 0) \text{ is eventually relayed by } (x,  h)}\right] = \e\left[\sum_{(x, h)\in \landsc} \mathds{1}_{(0, 0) \text{ is eventually relayed by } (x,  h)}\right] = \e[T].\]
Furthermore
\[\e\left[\sum_{(x, h)\in \tilde{\landsc}_1} \mathds{1}_{h = 0, (x, h) \text{ is eventually relayed by } (0,  H)}\right] = \e[\# \{x\in \Phi: (x, 0)\text{ is eventually relayed by } (0, H) \text{ in }\tilde{\landsc}\}].\]

Which finally leads to 
\[\lambda' \frac{\lambda}{\lambda+\lambda'} \e[\ell(\zone(0, H))] = \frac{\lambda'}{\lambda+\lambda'} \e[T].\]
Which is exactly the announced result.
\end{proof}

\end{document}
